\section{Triangulation for a manifold with a framed knot in it\\ and relative Pachner moves}
\label{sec_pachner}

We consider a closed oriented three-manifold~$M$ and a triangulation~$T$ of it containing a distinguished chain of two tetrahedra~$ABCD$ of one of the forms depicted in Figs.~\ref{fig1a} and~\ref{fig1b}. These two tetrahedra can either have the same orientation, as in Fig.~\ref{fig1a}, or the opposite orientations, as in Fig.~\ref{fig1b}.
\begin{figure}[ht]
\centerline{\includegraphics[scale=0.3]{figure1a}}
\caption{A chain of two identically oriented tetrahedra~$ABCD$.}
\label{fig1a}
\end{figure}

\begin{figure}[ht]
\centerline{\includegraphics[scale=0.3]{figure1b}}
\caption{A chain of two oppositely oriented tetrahedra~$ABCD$.}
\label{fig1b}
\end{figure}

Our construction of the invariant requires adopting the following convention (see Subsection~\ref{subs_daav} for details).

\begin{convention}
\label{conv1}
Any triangulation considered in this paper, including those which appear below at any step of a sequence of Pachner moves, is required to possess the following property: \emph{all} vertices of any tetrahedron are dif\/ferent.
\end{convention}

\begin{remark}
\label{remark-about-convention}
In particular, this convention is satisf\/ied by a \emph{combinatorial triangulation}, i.e., such a triangulation where any simplex is uniquely determined by the set of its vertices, all of those being dif\/ferent.
\end{remark}

To select a special chain of two tetrahedra as depicted in Figs.~\ref{fig1a} and~\ref{fig1b} essentially means the same as to select a \emph{framed knot} in~$M$. To be exact, there is a knot with \emph{two} framings given either by two closed lines (which we imagine as close to each other) $ACA$ and $DBD$, or by the two lines $ABA$ and $DCD$. In the case of the same orientation of the two tetrahedra, these possibilities lead to framings which dif\/fer in one full revolution (of the ribbon between two lines), so, to ensure the invariant character of choosing the framing, we have to choose the ``intermediate'' framing dif\/fering from them both in one-half of a revolution as the framing corresponding to our picture. In the case of the opposite orientations of the two tetrahedra, both ways simply give the same framing.

\begin{remark}
\label{remark-about-half-integers}
Thus, a \emph{half-integer} framing corresponds to Fig.~\ref{fig1a}, represented by a ribbon going like a M\"obius band, and an integer framing~-- to Fig.~\ref{fig1b}. These are the two kinds of framings we will be dealing with in this paper.
\end{remark}

Our aim is to construct an invariant of a pair $(M, K)$, where $K$ is a framed knot in~$M$, starting from a triangulation of~$M$ containing two distinguished tetrahedra as in Fig.~\ref{fig1a} or~\ref{fig1b}. To achieve this, we will construct in Section~\ref{sec_acyclic_general} a value not changing under Pachner moves on triangulation of~$M$ \emph{not touching the distinguished tetrahedra} of Fig.~\ref{fig1a} or~\ref{fig1b}. By ``not touching'' we understand those moves that do not replace either of the two tetrahedra in Fig.~\ref{fig1a} or~\ref{fig1b} with any other tetrahedra, and we call such moves \emph{relative Pachner moves}.

Recall that Pachner moves are elementary rebuildings of a \emph{closed} triangulated manifold. There are four such moves on three-dimensional manifolds. Two of them are illustrated in Figs.~\ref{fig2a} and~\ref{fig2b}, and the other two are inverse to these. The move in Fig.~\ref{fig2a} replaces the two adjacent tetrahedra $MNPQ$ and $RMNP$ with three new tetrahedra: $MNRQ$, $NPRQ$, and~$PMRQ$. The move in Fig.~\ref{fig2b} replaces one tetrahedron $MNPQ$ with four of them: $MNPR$, $MNRQ$, $MPQR$, and~$NPRQ$.
\begin{figure}[ht]
\centerline{\includegraphics[scale=0.39]{figure2a}}
\caption{A $2 \to 3$ Pachner move in three dimensions.}
\label{fig2a}
\end{figure}

\begin{figure}[ht]
\centerline{\includegraphics[scale=0.39]{figure2b}}
\caption{A $1 \to 4$ Pachner move in three dimensions.}
\label{fig2b}
\end{figure}

The main objective of the present section is to prove the following
\begin{theorem}
\label{th1}
Let M be a closed oriented three-manifold, $T_1$ and $T_2$ its triangulations with the same chain of two distinguished tetrahedra~$ABCD$, as depicted in Figs.~{\rm \ref{fig1a}} and~{\rm \ref{fig1b}}. Then, $T_1$~and~$T_2$ are related by a sequence of relative Pachner moves.
\end{theorem}

\begin{proof}
We are going to apply techniques from Lickorish's paper~\cite{L}. Therefore, let us f\/irst explain a method to subdivide the triangulations $T_1$ and $T_2$ in such way that they become combinatorial. Together with Pachner moves (see Figs.~\ref{fig2a} and~\ref{fig2b}), we will use \emph{stellar moves}, see~\cite[Section 3]{L}. In three dimensions, there is no dif\/f\/iculty to express the latter in terms of the former (and vice versa).

We can assume that $T_1$ and~$T_2$ already do not contain any more edges or two-dimensional faces whose vertices all lie in the set $\{A,B,C,D\}$ except those depicted in Figs.~\ref{fig1a} and~\ref{fig1b}~-- otherwise, we can always make obvious stellar subdivisions to ensure this.

Starting from the triangulation $T_1$, we f\/irst do Pachner moves $1\to 4$ in all tetrahedra adjacent to those two in Figs.~\ref{fig1a} and~\ref{fig1b}. Thus, there have appeared eight new vertices, we call them $N_1,\ldots,N_8$.

Next, we look at the edges in Figs.~\ref{fig1a} and~\ref{fig1b}. We are going to make some moves so that the link of each of them contain exactly one vertex between any two~$N_i$. If there were more such vertices, we can eliminate them from the link by doing suitable $2\to 3$ Pachner moves. Namely, the $2\to 3$ Pachner move provides a new edge joining $N_i$ directly with a ``farther'' vertex in the link, thus eliminating from the link the vertex next to $N_i$, see Fig.~\ref{fig3}.
A special case is two edges $AD$ and~$BC$: they require such procedure to be applied twice, ``on two sides''.
\begin{figure}[ht]
\centerline{\includegraphics[scale=0.29]{figure3}}
\caption{The edge drawn in boldface dashed line appears as a result of move $2\to 3$ and eliminates vertex~$P$ from the link of~$BD$.}
\label{fig3}\vspace{-2mm}
\end{figure}

This done, we make stellar subdivisions in the two-dimensional faces which are joins of the edges in Fig.~\ref{fig1a} or~\ref{fig1b} and the vertices lying between the $N_i$'s~-- two such vertices for each of $AD$ and~$BC$, and one for each of the remaining edges.

After that, we remove from the resulting simplicial complex those tetrahedra that have \emph{at least two} vertices in the set $\{A,B,C,D\}$. Specif\/ically, take these tetrahedra together with all their faces and denote by $L$ the obtained subcomplex. Then we remove~$L$ from our simplicial complex and take the closure of what remains; let~$V_1$ denote this closure.

Note that the triangulated boundary of~$V_1$ (which is also the boundary of~$L$) can be described as follows: f\/irst, \emph{double} the edges $AD$ and $BC$ in Figs.~\ref{fig1a} and~\ref{fig1b} in such way as to make a torus out of the boundary of the tetrahedron chain, and then make a barycentric subdivision of this triangulated torus.

Finally, we subdivide our simplicial complex~$V_1$, doing, e.g., suitable stellar moves in its simplices but leaving the boundary untouched, so that it becomes a combinatorial triangulation, as it is required in order to apply the techniques from~\cite{L}. Let~$W_1$ denote the resulting simplicial complex.

Now, we apply the above procedure to the triangulation $T_2$ and similarly obtain the simplicial complex $W_2$.
Obviously, $W_1$ and $W_2$ are PL-homeomorphic. Then, according to \cite[Theorem~4.5]{L}, these simplicial complexes are stellar equivalent. Moreover, there exists a chain of stellar and inverse stellar moves transforming $W_1$ into $W_2$ and such that \emph{the initial edges in the boundary~$\partial W_1$ never dissapear during the whole process}: they may be at most divided by several starrings done at them, but f\/inally the obtained fragments are glued together again; accordingly, the initial vertices in~$\partial W_1$~-- ends of initial edges~-- remain intact.

Indeed, the cited theorem in~\cite{L} is valid for simplicial complexes and not necessarily mani\-folds. Thus, we can do the following trick: glue to any edge in~$\partial W_1$ an additional two-cell~-- a triangle~-- by one of its sides, and do the same with edges in~$\partial W_2$. The obtained simplicial complexes $W'_1$ and~$W'_2$ are still~PL~-- and consequently stellar~-- homeomorphic, and obviously the additional cells are conserved (at most divided in parts but then glued again) along the whole chain of starrings and inverse starrings, marking thus the edges.

Such chain of starrings and inverse starrings can be extended to the union $W_1\cup L$, so that the result is $W_2\cup L$. Indeed, every stellar move involving~$\partial L$ is done either on one of initial edges or inside one of initial triangles. It is not hard to see that the extension from~$\partial L$ to~$L$ of a starring or inverse starring goes smoothly in both cases.

The subcomplex in Figs.~\ref{fig1a} and~\ref{fig1b} will not be touched by any move in the sequence. So, what remains is to replace the stellar moves with suitable (sequences of) Pachner moves.
\end{proof}


\section{Geometric values needed for the acyclic complex}
\label{sec_deficit_angles}

We are now going to construct an acyclic complex which produces the invariant of a three-manifold with a framed knot in it given by a chain of two tetrahedra as in Figs.~\ref{fig1a} and~\ref{fig1b}. The complex will be like those in \cite{TMP24, TMP15}, but in fact a bit simpler.

\begin{convention}
\label{conv2}
Recall that we are considering an \emph{orientable} manifold~$M$. From now on, we f\/ix a \emph{consistent orientation} for all tetrahedra in the triangulation. The orientation of a tetrahedron is understood here as an ordering of its vertices up to an even permutation; for instance, two tetrahedra $MNPQ$ and~$RMNP$, having a common face~$MNP$, are consistently oriented.
\end{convention}

\subsection{Oriented volumes and def\/icit angles}
\label{subs_daav}
We need the so-called \emph{deficit angles} corresponding to the edges of triangulation. The rest of this section is devoted to explaining these def\/icit angles and related notions, while the acyclic complex itself will be presented in Section~\ref{sec_acyclic_general}.

Recall that we assume that \emph{all} the vertices of any tetrahedron in the triangulation are dif\/ferent (see Convention~\ref{conv1}). Put all the vertices of the triangulation in the Euclidean space $\mathbb R^3$ (i.e., we assign to each of them three real coordinates) in an arbitrary way with only one condition: the vertices of each tetrahedron in the triangulation must not lie in the same plane. This condition ensures that the geometric quantities we will need~-- edge lengths and tetrahedron volumes~-- never vanish.

When we put an oriented tetrahedron $MNRQ$ into~$\mathbb R^3$, we can ascribe to it an \emph{oriented volume} denoted $V_{MNRQ}$ according to the formula
\begin{equation}
6V_{MNRQ}=\overrightarrow{MN} \cdot \overrightarrow{MQ} \cdot \overrightarrow{MR}
\label{eq_oriented_volume}
\end{equation}
(scalar triple product in the right-hand side). If the sign of the volume def\/ined by~(\ref{eq_oriented_volume}) of a given tetrahedron is positive, we say that it is put in~$\mathbb R^3$ \emph{with its orientation preserved}; if it is negative we say that it is put in~$\mathbb R^3$ \emph{with its orientation changed}.

Now we consider the dihedral angles at the edges of triangulation. We will ascribe a sign to each of these angles \emph{coinciding with
the oriented volume sign} of the tetrahedron to which the angle belongs. Consider a certain edge $QR$ in the triangulation, and let its link contain vertices $P_1, \ldots, P_n$, so that the tetrahedra $P_1P_2RQ, \ldots, P_nP_1RQ$ are situated around~$QR$ and form its \emph{star}. With our def\/inition for the signs of dihedral angles, one can observe that the algebraic sum of all angles at the edge~$QR$ is a multiple of $2\pi$, if these angles are calculated according to the usual formulas of Euclidean geometry, starting from \emph{given coordinates of vertices} $P_1, \ldots, P_n$, $Q$ and~$R$.

Namely, we would like to use the following method for computing these dihedral angles. Given the coordinates of vertices, we calculate all the \emph{edge lengths} in tetrahedra $P_1P_2RQ, \ldots$, $P_nP_1RQ$ and the signs of all tetrahedron volumes, and then we calculate dihedral angles from the edge lengths.

Suppose now that we have slightly, but otherwise arbitrarily, changed the edge lengths. Each separate tetrahedron $P_1P_2RQ, \ldots, P_nP_1RQ$ remains still a Euclidean tetrahedron, but the algebraic sum of their dihedral angles at the edge~$QR$ ceases, generally speaking, to be a~multiple of $2\pi$. This means that these tetrahedra can no longer be put in~$\mathbb R^3$ together. In such situation, we call this algebraic sum, taken with the opposite sign, \emph{deficit (or defect) angle} at edge~$QR$:
\begin{equation}
\omega_{QR} = - \sum_{i=1}^n \varphi_i \bmod 2\pi,
\label{eq_deficit_angle}
\end{equation}
where $\varphi_i$ are the dihedral angles at $QR$ in the $n$ tetrahedra under consideration. Note that the minus sign in~(\ref{eq_deficit_angle}) is just due to a convention in ``Regge calculus''~\cite{Regge} where such def\/icit angles often appear.


\subsection{A relation for inf\/initesimal deformations of def\/icit angles}

To build our chain complex~(\ref{eq_alg_complex}) in Section~\ref{sec_acyclic_general}, we need only \emph{infinitesimal} def\/icit angles arising from inf\/initesimal deformations of edge lengths in the neighborhood of a f\/lat case, where all def\/icit angles vanish.

\begin{lemma}
\label{lemma_infinitesimal}
Let $Q$ be any vertex in a triangulation of closed oriented three-manifold, and $QR_1,\ldots,\allowbreak QR_m$ all the edges that end in~$Q$. If infinitesimal deficit angles $d\omega_{QR_i}$ are obtained from infinitesimal deformations of length of edges in the triangulation with respect to the flat case, then{\samepage
\begin{equation}
\sum_{i=1}^m \vec e_{QR_i}\, d\omega_{QR_i} = \vec 0,
\label{eq_zero_for_dx*}
\end{equation}
where $\vec e_{QR_i}=1/l_{QR_i} \cdot \overrightarrow{QR_i}$ is a unit vector along $\overrightarrow{QR_i}$.}
\end{lemma}

\begin{proof}
We f\/irst consider the case of small but f\/inite deformations of edge lengths and def\/icit angles.

Let the edge lengths in tetrahedra $P_1P_2R_iQ, \ldots, P_nP_1R_iQ$, forming the star of edge~$QR_i$, be slightly deformed with respect to the f\/lat case $\omega_{QR_i}=0$. We can introduce a Euclidean coordinate system in tetrahedron~$P_1P_2R_iQ$. Then, this coordinate system can be extended to tetrahedron~$P_2P_3R_iQ$ through their common face~$P_2R_iQ$. Continuing this way, we can go around the edge~$QR_i$ and return in the initial tetrahedron~$P_1P_2R_iQ$, obtaining thus a new coordinate system in it. The transformation from the old system to the new one is an orthogonal rotation around the edge~$QR_i$ through the angle~$\omega_{QR_i}$ in a proper direction.

More generally, we can consider going around edge~$QR_i$ but starting from some arbitrary ``remote'' tetrahedron~$\mathcal T \subset \St(Q)$ and going f\/irst from~$\mathcal T$ to~$P_1P_2R_iQ$ through some two-dimen\-sion\-al faces in the triangulation, then going around~$QR_i$ as in the preceding paragraph, and f\/inally returning from~$P_1P_2R_iQ$ to~$\mathcal T$ following the same way. The corresponding transformation from the old coordinate system to the new one will be again an orthogonal rotation around some axis going through~$Q$.

Now, imagine we are within some chosen tetrahedron~$\mathcal T$ belonging to the star of~$Q$ (i.e., having $Q$ as one of its vertices). For clarity, let us be at a point~$N$ on a small sphere around~$Q$. We are going to draw some special closed paths on this sphere minus $m$~punctures~-- the points where the sphere intersects the edges~$QR_i$. It is clear that we can choose a set of paths $\alpha_i$, $i=1,\dots,m$, in such punctured sphere, with the following properties:
\begin{itemize}
\itemsep=0pt
	\item each $\alpha_i$ begins and ends in~$N$,
	\item within each $\alpha_i$ lies exactly one puncture, namely the intersection of sphere with~$QR_i$, and $\alpha_i$ goes around it in the counterclockwise direction,
	\item the composition $\alpha_1 \circ \dots \circ \alpha_m$ is homotopic to the trivial path in the punctured sphere staying at the point~$N$ .	
\end{itemize}

Denote $X_{QR_i}$ the coordinate system transformation corresponding to~$\alpha_i$. Choosing vertex~$Q$ as the origin of coordinates, we can represent all $X_{QR_1}, \ldots, X_{QR_m}$ as elements of the group~${\rm SO}(3)$. The collapsibility of the composition of~$\alpha_i$ leads to the relation
\begin{equation}
X_{QR_m} \circ \dots \circ X_{QR_1} = \mathop{\rm Id}
\label{eq_rel_in_SO3}
\end{equation}
(the inverse order of $X_{QR_i}$ compared to~$\alpha_i$ ref\/lects the fact that, in the product in~\eqref{eq_rel_in_SO3}, the rightmost elements comes f\/irst, while in the product of~$\alpha_i$~-- traditionally the leftmost).

Now we turn to the inf\/initesimal case. Here, up to second-order inf\/initesimals, $X_{QR_i}$ are just rotations around corresponding edges~$QR_i$ (regardless of how exactly the path~$\alpha_i$ goes between the other edges):
\begin{displaymath}
X_{QR_i} = \mathop{\rm Id} + x_{QR_i}\, d\omega_{QR_i},
\end{displaymath}
where $x_{QR_i}$ is the element of Lie algebra $\mathfrak{so}(3)$ def\/ined by
\begin{displaymath}
(x_{QR_i})_{\alpha\beta} = \sum\limits_\gamma\varepsilon_{\alpha\beta\gamma}\, (\vec e_{QR_i})_\gamma ,
\end{displaymath}
$\varepsilon_{\alpha\beta\gamma}$ being the Levi-Civita symbol.
Equation~(\ref{eq_rel_in_SO3}) gives
\begin{displaymath}
\sum_{i=1}^m x_{QR_i}\, d\omega_{QR_i} = 0 ,
\end{displaymath}
so we get~(\ref{eq_zero_for_dx*}).
\end{proof}


\subsection[Formulas for $\partial \omega_a/\partial l_b$]{Formulas for $\boldsymbol{\partial \omega_a/\partial l_b}$}

The main ingredient of our complex~(\ref{eq_alg_complex}) is matrix $f_3 = \big(\frac{\partial \omega_a}{\partial l_b}\big)$, where $a$ and $b$ run over all edges in the triangulation and the partial derivatives are taken at such values of all lengths~$l_b$ where all $\omega_a = 0$. We are going to express these derivatives in terms of edge lengths and tetrahedron volumes. Recall that all the tetrahedra we deal with are supposed to be consistently oriented (see Convention~\ref{conv2}), so that their volumes have signs. Nonzero derivatives are obtained in the following cases.

\subsubsection*{1st case}
Edges $a=MN$ and $b=QR$ are skew edges of tetrahedron $MNRQ$:
\begin{equation}
\frac{\partial \omega_{MN}}{\partial l_{QR}} = -\frac {l_{MN}\,l_{QR}}{6}\, \frac {1}{V_{MNRQ}}.
\label{eq MNQR}
\end{equation}
This formula is an easy exercise in elementary geometry.

\subsubsection*{2nd case}
Edges $a=MN$ and $b=MP$ belong to the common face of two neighboring tetrahedra $MNPQ$ and $RMNP$ (see the left-hand side of Fig.~\ref{fig2a}):
\begin{equation}
\frac{\partial \omega_{MN}}{\partial l_{MP}} = \frac {l_{MN}\,l_{MP}}{6}\,\frac {V_{NPRQ}}{V_{MNPQ}\,V_{RMNP}}.
\label{eq MNMP}
\end{equation}
To prove this formula, suppose the lengths $l_{MP}$ and $l_{QP}$ are free to change, while the lengths of the remaining seven edges in the two mentioned tetrahedra, \emph{and also the length~$l_{QR}$}, are f\/ixed (this implies of course f\/ixing the dihedral angles at edge~$MN$). Then $l_{QP}$ is a function of~$l_{MP}$ and, using formulas of type~(\ref{eq MNQR}), one can f\/ind that
\begin{equation}
\frac{\partial l_{QP}}{\partial l_{MP}} = \frac{l_{MP}}{l_{QP}}\, \frac{V_{NPRQ}}{V_{RMNP}}.
\label{eq ll1}
\end{equation}
Then,
\begin{equation*}
0 \equiv \frac{d\omega_{MN}}{dl_{MP}} = \frac{\partial \omega_{MN}}{\partial l_{MP}} + \frac{\partial \omega_{MN}}{\partial l_{QP}}\, \frac{\partial l_{QP}}{\partial l_{MP}} = \frac{\partial \omega_{MN}}{\partial l_{MP}} - \frac {l_{MN}\,l_{MP}}{6}\,\frac {V_{NPRQ}}{V_{MNPQ}\,V_{RMNP}}
\end{equation*}
and formula~(\ref{eq MNMP}) follows.

\subsubsection*{3rd case}
Edge $a=b=QR$ is common for exactly three tetrahedra $MNRQ$, $NPRQ$ and~$PMRQ$ (as in the right-hand side of Fig.~\ref{fig2a}):
\begin{equation}
\frac{\partial \omega_{QR}}{\partial l_{QR}} = -\frac {l_{QR}^2}{6}\, \frac {V_{MNPQ}\,V_{RMNP}}{V_{MNRQ}\,V_{NPRQ}\,V_{PMRQ}}.
\label{eq QRQR}
\end{equation}
Observe that this is the most important formula allowing us to construct manifold invariants based on three-dimensional Euclidean geometry. To prove it, consider the ten edges in the three mentioned tetrahedra and suppose that the lengths $l_{MP}$ and $l_{QR}$ are free to change, while the lengths of the remaining eight edges are f\/ixed. Then,
\begin{equation}
\frac{\partial l_{MP}}{\partial l_{QR}} = -\frac{l_{QR}}{l_{MP}}\, \frac{V_{MNPQ}\,V_{RMNP}}{V_{MNRQ}\, V_{NPRQ}}.
\label{eq ll0}
\end{equation}

To prove~\eqref{eq ll0}, we use the following trick: f\/irst, consider the situation where $l_{QP}$ and~$l_{MP}$ are free to change, and the remaining eight lengths are f\/ixed. Then, $\partial l_{QP} / \partial l_{MP}$ is given by formula~\eqref{eq ll1}. Second, consider the situation where $l_{QP}$ and~$l_{QR}$ are free to change, and the remaining eight lengths are f\/ixed. Now a formula of type~\eqref{eq ll1} gives
\begin{equation}
\frac{\partial l_{QP}}{\partial l_{QR}} = \frac{l_{QR}}{l_{QP}}\, \frac{V_{NQMP}}{V_{MNRQ}}.
\label{eq n}
\end{equation}
Third, let now the \emph{seven} edge lengths except $l_{QP}$, $l_{MP}$ and~$l_{QR}$ be f\/ixed; consider $l_{QP}$ as a~function of the other two and then equate $l_{QP}$ to a constant:
\begin{equation*}
l_{QP}(l_{MP},l_{QR}) = \const .
\end{equation*}
This can be viewed as def\/ining $l_{MP}$ as an \emph{implicit function} of~$l_{QR}$, and its derivative, calculated according to the standard formula and using \eqref{eq ll1} and~\eqref{eq n}, is exactly~\eqref{eq ll0}.

Now we can write (assuming again that only $l_{MP}$ and $l_{QR}$ can change):
\begin{equation*}
0 \equiv \frac{d\omega_{QR}}{dl_{QR}} = \frac{\partial \omega_{QR}}{\partial l_{QR}} + \frac{\partial \omega_{QR}}{\partial l_{MP}}\, \frac{\partial l_{MP}}{\partial l_{QR}} = \frac{\partial \omega_{QR}}{\partial l_{QR}} + \frac {l_{QR}^2}{6}\, \frac{V_{MNPQ}\, V_{RMNP}}{V_{MNRQ}\, V_{NPRQ}\, V_{PMRQ}} \, ,
\end{equation*}
and formula~(\ref{eq QRQR}) follows.

\subsubsection*{4th case}
Edge $a=b=QR$ is common for $n > 3$ tetrahedra $P_{i-1}P_iRQ$ $(i = 1,\ldots,n)$:
\begin{equation}
\frac{\partial \omega_{QR}}{\partial l_{QR}} = -\frac {l_{QR}^2}{6}\, \sum\limits_{i=3}^n
\frac {V_{P_1P_{i-1}P_iQ}\,V_{RP_1P_{i-1}P_i}}{V_{P_1P_{i-1}RQ}\,V_{P_{i-1}P_iRQ}\,V_{P_iP_1RQ}}.
\label{eq QRQRn}
\end{equation}
This formula comes out if we draw $n-3$ diagonals $P_1P_i$ ($i = 3,\ldots,n-1$) and apply~(\ref{eq QRQR}) to each of f\/igures $P_1P_{i-1}P_iQR$. Adding up the def\/icit angles around $QR$ in each of those f\/igures and cancelling out the dihedral angles which enter twice and with the opposite signs, we get nothing else than $\omega_{QR}$ for the whole f\/igure.


\section{The acyclic complex and the invariant}
\label{sec_acyclic_general}

\subsection{Generalities on acyclic complexes and their Reidemeister torsions}

We brief\/ly review basic def\/initions from the theory of based chain complexes and their associated Reidemeister torsions, see monograph~\cite{Turaev:2001} for details.

\begin{definition}
\label{dfn:cmplx}
Let $C_0, C_1, \ldots , C_n$ be f\/inite-dimensional $\mathbb R$-vector spaces. The sequence of vector spaces and linear mappings
\begin{equation}
\label{eq:cmplx}
C_* = 0 \xrightarrow{f_n} C_n \xrightarrow{f_{n-1}} C_{n-1} \xrightarrow{} \cdots \xrightarrow{} C_1 \xrightarrow{f_0} C_0 \xrightarrow{f_{-1}} 0
\end{equation}
is called a \emph{chain complex} if $\Ima f_i \subset \Ker f_{i-1}$ for all $i = 0, \ldots, n$. This condition is equivalent to $f_{i-1} \circ f_i = 0$.
\end{definition}

Suppose that chain complex $C_*$ is \emph{based}, i.e., each $C_i$ is endowed with a distinguished basis. To make the notations in this paper consistent with our previous papers, we denote this basis~$\mathcal C_i$ (rather than~$\mathbf c_i$, which is the usual notation in the literature on torsions). The linear mapping $f_i \colon C_{i + 1} \to C_i$ can thus be identif\/ied with its matrix.

\begin{definition}
\label{dfn:acycl}
The based chain complex $C_*$ is said to be \emph{acyclic} if $\Ima f_i = \Ker f_{i-1}$ for all~$i = 0, \ldots, n$.
\end{definition}

\begin{remark}
This condition is equivalent to $\rank f_{i-1} = \dim C_i - \rank f_i$. Since $\Ima f_n = 0$ and $\Ker f_{-1} = C_0$, it follows that in an acyclic complex $f_{n-1}$ is injective and~$f_0$ is surjective.
\end{remark}

Suppose that the chain complex $C_*$ def\/ined in~\eqref{eq:cmplx} is acyclic. For every $i = 0, \ldots, n$, let $\mathcal B_i$ be a subset of the basis~$\mathcal C_i$ such that $f_{i-1} (\mathcal B_i)$ is a basis for~$\Ima f_{i-1}$. Recall that we consider the linear operator~$f_i$ as a matrix whose columns and rows correspond to the basis vectors in~$C_{i+1}$ and~$C_i$ respectively. Denote by $f_i|_{\mathcal B_{i+1}, \overline{\mathcal B}_i}$ the submatrix of~$f_i$ consisting of columns corresponding to $\mathcal B_{i+1}$ and rows corresponding to~$\overline{\mathcal B}_i = \mathcal C_i \setminus \mathcal B_i$.

Due to the acyclicity, $\#\mathcal B_{i+1} = \rank f_i = \dim C_i - \rank f_{i-1} = \dim C_i - \#\mathcal B_i$ and it follows that $f_i|_{\mathcal B_{i+1}, \overline{\mathcal B}_i}$ is square. It is also nondegenerate: its determinant coincides with the determinant of the transition matrix between the two bases $\mathcal C_i$ and $\mathcal B_i \cup f_i(\mathcal B_{i+1})$ for~$C_i$.

\begin{definition}
\label{dfn:tors}
The sign-less quantity
\begin{equation}
\label{eq:torsion}
\tau(C_*) = \prod\limits_{i = 0}^{n-1} \bigl(\det f_i|_{\mathcal B_{i+1}, \overline{\mathcal B}_i}\bigr)^{(- 1)^{i+1}} \in \mathbb{R}^* / \{\pm 1\}
\end{equation}
is called the \emph{Reidemeister torsion} of the acyclic based chain complex~$C_*$.
\end{definition}

\begin{remark}[\cite{Turaev:2001}]
It is easy to show that $\tau(C_*)$ does not depend on the choice of the subsets~$\mathcal B_i$, but of course depends on the choice of the bases in~$C_i$'s.
\end{remark}

\begin{remark}\label{rem1B}
We will also use simplif\/ied notations for~$f_i|_{\mathcal B_{i+1}, \overline{\mathcal B}_i}$, explaining their meaning in the text, as in equation~\eqref{eq_tau} below.
\end{remark}

\subsection{The acyclic complex}

Consider the following sequence of vector spaces and linear mappings:
\begin{equation}
0\xrightarrow{} (dx)'\xrightarrow{f_2} (dl)' \xrightarrow{f_3=f_3^{\rm T}}(d\omega)' \xrightarrow{f_4=-f_2^{\rm T}} \oplus' \mathfrak{so}(3) \xrightarrow{} 0\,.
\label{eq_alg_complex}
\end{equation}
Here is the detailed description of the vector spaces in the chain complex~(\ref{eq_alg_complex}):
\begin{itemize}
\itemsep=0pt
  \item the f\/irst vector space $(dx)'$ is the vector space spanned by the dif\/ferentials of coordinates of all vertices except $A$, $B$, $C$ and~$D$;
  \item the second vector space $(dl)'$ is the vector space spanned by the dif\/ferentials of edge lengths for all edges except those depicted in Figs.~\ref{fig1a} and~\ref{fig1b};
  \item similarly, the third vector space $(d\omega)'$ is the vector space spanned by the dif\/ferentials of def\/icit angles corresponding to the same edges;
  \item the last vector space is a direct sum of copies of the Lie algebra $\mathfrak{so}(3)$ corresponding to the same vertices in the triangulation as in the f\/irst space.
\end{itemize}

Before giving the detailed def\/initions of mappings $f_2$, $f_3$ and~$f_4$, here are some remarks.
\begin{remark}
We use the notations ``$f_2$'' and ``$f_3$'' to make them consistent with other papers on the subject, such as, e.g.~\cite{M,M-thesis}. So, the reader must not be surprised with not f\/inding any~``$f_1$'' in this paper.
\end{remark}
\begin{remark}
\label{rem-b}
There is a natural basis in each of the vector spaces, given by the corresponding dif\/ferentials in $(dx)'$, $(dl)'$ and~$(d\omega)'$, and by the standard Lie algebra generators in~$\oplus' \mathfrak{so}(3)$. Such basis is determined up to an ordering of the vertices in $(dx)'$ and~$\oplus' \mathfrak{so}(3)$, and up to an ordering of the edges in $(dl)'$ and~$(d\omega)'$. Thus, we can identify the elements of vector spaces with column vectors, and mappings~-- with matrices.

For example, the vector space~$(dx)'$ consists of columns of the kind
\begin{displaymath}
\big(dx_{E_1}, dy_{E_1}, dz_{E_1}, \ldots, dx_{E_{N'_0}}, dy_{E_{N'_0}}, dz_{E_{N'_0}}\big)^{\rm T},
\end{displaymath}
where $E_1,\ldots,E_{N'_0}$ are all the vertices in the triangulation except $A$, $B$, $C$ and~$D$. We use thus notation~$N'_0$ for the number of vertices that are \emph{inner} for the manifold with boundary ``$M$~minus the interior of two tetrahedra~$ABCD$ (and with doubled edges $AD$ and~$BC$)'', which is consistent with our other papers, where simply~$N_0$ denotes the number of \emph{all} vertices. Note also that the special role of the edges depicted in Figs.~\ref{fig1a} and~\ref{fig1b}~-- boundary edges for the mentioned manifold, if we double $AD$ and~$BC$~-- consists in the fact that they simply do not take part in forming the second and third linear spaces in complex~(\ref{eq_alg_complex}).
\end{remark}
\begin{remark}
The superscript~$\rm T$ means matrix transposing; the equalities over the arrows in complex~(\ref{eq_alg_complex}) will be proved soon after we def\/ine the mappings $f_2$, $f_3$ and~$f_4$, see Theorem~\ref{th_symmetry}.
\end{remark}
\begin{remark}
The fact that the $f$'s are numbered in the decreasing order in~\eqref{eq:cmplx} and in the increasing order in~\eqref{eq_alg_complex} brings no dif\/f\/iculties.
\end{remark}
\begin{remark}
In our case, there will be no problems with the sign of the torsion, and thus the factoring by~$\{ \pm 1 \}$, as in~\eqref{eq:torsion}, will be unnecessary, see Subsection~\ref{subsec-torsion-invariant} for details.
\end{remark}

Here are the def\/initions of the mappings in the chain complex~(\ref{eq_alg_complex}):
\begin{itemize}
\itemsep=0pt
\item the def\/inition of mapping $f_2$ is obvious: if we change inf\/initesimally the coordinates of vertices, then the corresponding edge length changes are obtained by dif\/ferentiating formulas of the kind
\begin{equation}
l_{MN} = \sqrt{(x_N-x_M)^2+(y_N-y_M)^2+(z_N-z_M)^2},
\label{eq_1a}
\end{equation}
where $M$ and $N$ are two vertices, $x_M,\ldots,z_N$~-- their coordinates, and $l_{MN}$~-- the length of edge~$MN$;
\item mapping $f_3$ goes according to formulas~(\ref{eq MNQR})--(\ref{eq QRQRn});
\item for the mapping~$f_4$, the element of the Lie algebra corresponding to a given vertex, arising from given curvatures~$d\omega$ due to~$f_4$, is by def\/inition given by the left-hand side of formula~(\ref{eq_zero_for_dx*}).
\end{itemize}


\begin{theorem}
Sequence \eqref{eq_alg_complex} is a chain complex, i.e., the composition of two successive mappings is zero.
\label{th2}
\end{theorem}

\begin{proof}
The equality $f_3 \circ f_2 = 0$ is obvious from geometric considerations. Indeed, the edge length changes caused by changes of vertex coordinates give no def\/icit angles, because the whole picture (vertices and edges) does not go out of the Euclidean space~$\mathbb R^3$.

By Lemma~\ref{lemma_infinitesimal}, the equality $f_4 \circ f_3 = 0$ holds as well.
\end{proof}

Complex~(\ref{eq_alg_complex}) can be called a \emph{complex of infinitesimal geometric deformations}. It turns out to have an interesting symmetry property.
\begin{theorem}\label{th_symmetry}
The matrices of mappings in complex \eqref{eq_alg_complex} satisfy the following symmetry properties:
\begin{equation}
f_3 = f_3^{\rm T},\qquad f_4 = - f_2^{\rm T}.
\label{eq_symmetry}
\end{equation}
\end{theorem}

\begin{proof} The symmetry of matrix~$f_3$ means that $\partial \omega_b / \partial l_c = \partial \omega_c / \partial l_b$, and it follows directly from explicit formulas~\eqref{eq MNQR} and \eqref{eq MNMP}.

As for the second equality in~(\ref{eq_symmetry}), it can be proved by a direct writing out of matrix elements, i.e., the relevant partial derivatives. For the mapping~$f_2$, one has to dif\/ferentiate relation~(\ref{eq_1a}), and for~$f_4$~-- use the left-hand side of~(\ref{eq_zero_for_dx*}).
\end{proof}

\begin{remark}
There exists also a dif\/ferent proof of the equality $f_3 = f_3^{\rm T}$, which enables us to look at it perhaps from a dif\/ferent perspective, and based on the \emph{Schl\"afli differential identity} for a Euclidean tetrahedron:
\begin{equation}
\sum_{i=1}^6 l_i \, d\varphi_i = 0
\label{eq_Schlaefli_0}
\end{equation}
for any inf\/initesimal deformations ($l_i$ are edge lengths in the tetrahedron, and~$\varphi_i$ are dihedral angles at edges). It follows from~\eqref{eq_Schlaefli_0} that
\begin{equation}
\sum_a l_a \, d\omega_a = 0,
\label{eq_Schlaefli_1}
\end{equation}
where $a$ runs over all edges in the triangulation.

Consider now the quantity $\Phi=\sum_a l_a \omega_a$ as a function of the lengths~$l_a$, and write the following identity for it:
\begin{equation}
\frac{\partial^2 \Phi}{\partial l_b \partial l_c} = \frac{\partial^2 \Phi}{\partial l_c \partial l_b},
\label{eq_Schlaefli_2}
\end{equation}
where $b$ and $c$ are some edges. It is easy to see that~(\ref{eq_Schlaefli_1}) together with~(\ref{eq_Schlaefli_2}) yield exactly the required symmetry.
\end{remark}

The examples below in Sections~\ref{sec_unknot} and~\ref{sec_lenses} show that there are many enough interesting cases where complex~(\ref{eq_alg_complex}) turns out to be acyclic (see Def\/inition~\ref{dfn:acycl}). However, at this time we cannot make this statement more precise.

\begin{convention}
\label{conv3}
From now on, we assume that we are working with an \emph{acyclic complex}.
\end{convention}


\subsection{The Reidemeister torsion and the invariant}
\label{subsec-torsion-invariant}

As complex~(\ref{eq_alg_complex}) is supposed to be acyclic, we associate to it its \emph{Reidemeister torsion} given by
\begin{equation}
\tau = \frac{\bigl(\det f_2 |_{\overline{\mathcal B}}\bigr) \bigl(\det f_4 |_{\overline{\mathcal B}}\bigr)}{\det f_3 |_{\mathcal B}} = (-1)^{N'_0}\, \frac{\bigl(\det f_2 |_{\overline{\mathcal B}}\bigr)^2}{\det f_3 |_{\mathcal B}}
\label{eq_tau}
\end{equation}
(cf. Def\/inition~\ref{dfn:tors} and Remark~\ref{rem1B} after it). Recall that $N'_0$, according to Remark~\ref{rem-b}, is the number of vertices in the triangulation without $A,B,C$ and~$D$. The letter $\mathcal B$ denotes a maximal subset of edges (and thus basis vectors in $(dl)'$ and~$(d\omega)'$; remember that the edges depicted in Figs.~\ref{fig1a} and~\ref{fig1b} have been already withdrawn) for which the corresponding \emph{ diagonal} minor of~$f_3$ does not vanish. We write this minor as $\det f_3|_{\mathcal B}$, where $f_3|_{\mathcal B}$ is the submatrix of $f_3$ whose \emph{rows and columns} correspond to the edges in~$\mathcal B$. The set $\overline{\mathcal B}$ is the complement of~$\mathcal B$ in the set of all edges except those depicted in Figs.~\ref{fig1a} and~\ref{fig1b}, and $f_2|_{\overline{\mathcal B}}$ (resp. $f_4|_{\overline{\mathcal B}}$) is the submatrix of~$f_2$ (resp.~$f_4$) whose \emph{rows} (resp.\ \emph{columns}) correspond to the edges in~$\overline{\mathcal B}$.

\begin{remark}
\label{remark_sign}
As it is known~\cite{Turaev:2001}, usually the Reidemeister torsion is def\/ined up to a sign, so that special measures must be taken for its ``sign-ref\/ining''. This sign is changed when we change the order of basis vectors in any of the vector spaces. In our case, however, this is not a problem: due to the symmetry proved in Theorem~\ref{th_symmetry}, we can choose our torsion in the form~(\ref{eq_tau}) where the numerator is a square and the denominator is a \emph{diagonal} minor. Both thus do not depend on the order of basis vectors.
\end{remark}

We now def\/ine the value
\begin{equation}
I(M) = \tau\, \frac{\prod_{\rm edges}' l^2}{\prod_{\rm tetrahedra}' (-6V)}\, (6V_{ABCD})^4.
\label{eq_invariant}
\end{equation}
Here $\prod_{\rm edges}' l^2$ means the product of squared lengths for all edges except those depicted in Figs.~\ref{fig1a} and~\ref{fig1b}, or simply \emph{inner} edges; $\prod_{\rm tetrahedra}' (-6V)$ means the product of all tetrahedron volumes multiplied by~$(-6)$ except two distinguished tetrahedra~$ABCD$, and $\tau$ is the Reidemeister torsion of complex (\ref{eq_alg_complex}) given by~(\ref{eq_tau}).

\begin{theorem}
Let $T_1$ and $T_2$ be triangulations of the manifold~$M$ with the same chain of two distinguished tetrahedra~$ABCD$ in them. If complex~\eqref{eq_alg_complex} is acyclic for $T_1$, then it is also acyclic for $T_2$, and the value $I(M)$ given by~\eqref{eq_invariant} is the same for $T_1$ and~$T_2$.
\label{th3}
\end{theorem}

\begin{proof}
By Theorem~\ref{th1}, it is enough to show the invariance of~$I(M)$ under relative Pachner moves.

Suppose we are doing a $2\to 3$ Pachner move: two adjacent tetrahedra $MNPQ$ and~$RMNP$ are replaced by three tetrahedra $MNRQ$, $NPRQ$
and $PMRQ$, see Fig.~\ref{fig2a}. Thus, a new edge~$QR$ appears in the triangulation. We are going to express the new matrix~$f_3$ in terms of the old one. Essentially, we follow~\cite[Section~4]{JNMP1}.

Set $\tilde l_{QR} = l_{QR} - l^{(0)}_{QR}$, where $l^{(0)}_{QR}$ is the solution of $\omega_{QR} = 0$, considered as a function of other edge lengths. Then,
\begin{equation}
d\tilde l_{QR} = dl_{QR} - \sum\limits_{i=1}^{N'_1} \frac{\partial l_{QR}}{\partial l_i}{\Big|}_{\omega_{QR} = 0} \, dl_i,
\label{eq_u}
\end{equation}
where $N'_1$ is the number of inner edges (we reserve the notation~$N_1$ for the number of \emph{all} edges, in analogy with our notations $N'_0$ and~$N_0$ for vertices) in the triangulation before doing the move $2\to 3$. As $\tilde l_{QR}=0$ implies $\omega_{QR}=0$, the dif\/ferential $d\omega_{QR}$ depends only on $d\tilde l_{QR}$ and does not depend on the rest of~$dl_i$. This yields
\begin{equation}
\begin{pmatrix}
d\omega_1 \\
\vdots \\
d\omega_{N'_1} \\
d\omega_{QR}
\end{pmatrix} =
\begin{pmatrix}
{} & {} & {} & \frac{\partial\omega_1}{\partial l_{QR}} \\
{} & f_3^{\rm old} & {} & \vdots \\
{} & {} & {} & \frac{\partial\omega_{N'_1}}{\partial l_{QR}} \\
0 & \cdots & 0 & \frac{\partial\omega_{QR}}{\partial l_{QR}}
\end{pmatrix}
\begin{pmatrix}
dl_1 \\
\vdots \\
dl_{N'_1} \\
d\tilde l_{QR}
\end{pmatrix},
\label{eq_v}
\end{equation}
where $f_3^{\rm old}$ is, of course, the matrix~$f_3$ before the move. Formulas \eqref{eq_v} and~\eqref{eq_u} give the representation of the ``new'' matrix $f_3^{\rm new}$ as the following product:
\begin{equation*}
f_3^{\rm new} =
\begin{pmatrix}
{} & {} & {} & \frac{\partial\omega_1}{\partial l_{QR}} \\
{} & f_3^{\rm old} & {} & \vdots \\
{} & {} & {} & \frac{\partial\omega_{N'_1}}{\partial l_{QR}} \\
0 & \cdots & 0 & \frac{\partial\omega_{QR}}{\partial l_{QR}}
\end{pmatrix}
\begin{pmatrix}
1 & {} & {} & 0 \\
{} & \ddots & {} & \vdots \\
{} & {} & 1 & 0 \\
- \frac{\partial l_{QR}}{\partial l_1} & \cdots & - \frac{\partial l_{QR}}{\partial l_{N'_1}} & 1
\end{pmatrix}.
\end{equation*}

The ``new'' set of edges~$\mathcal B^{\rm new}$ in~\eqref{eq_tau} can be taken as $\mathcal B^{\rm new} = \mathcal B^{\rm old} \cup\{QR\}$, then the minors of $f_2$ and~$f_4$ remain the same. Using~\eqref{eq QRQR}, we obtain the ratio between the new and old minors of~$f_3$:
\begin{equation}
\frac{\det f_3^{\rm new} |_{\mathcal B^{\rm new}}}{\det f_3^{\rm old} |_{\mathcal B^{\rm old}}} = \frac{\partial\omega_{QR}}{\partial l_{QR}} = - \frac{l_{QR}^2}{6}\frac{V_{MNPQ}V_{RMNP}}{V_{MNRQ}V_{NPRQ}V_{PMRQ}}  .
\label{eq QRQR minors}
\end{equation}
This together with \eqref{eq_invariant} and~\eqref{eq_tau} proves that $I(M)$ does not change under $2\to 3$ and $3\to 2$ Pachner moves.

Now we consider a $1\to 4$ Pachner move. It means that a tetrahedron $MNPQ$ is divided into four tetrahedra by adding a new vertex~$R$ inside it, as in Fig.~\ref{fig2b}. Hence, three new components are added to the vectors in the space~$(dx)'$~-- the f\/irst vector space in sequence~(\ref{eq_alg_complex}), namely the dif\/ferentials $dx_R$, $dy_R$ and~$dz_R$ of coordinates of vertex~$R$. In the same way, four new components are added to the vectors in the space~$(dl)$~-- the second vector space in~(\ref{eq_alg_complex}), namely $dl_{MR}$, $dl_{NR}$, $dl_{PR}$ and~$dl_{QR}$. We add the edge $QR$ to the set~$\mathcal B$, then $MR$, $NR$ and~$PR$ are added to~$\overline{\mathcal B}$. The minor of $f_2$ acquires a block triangular structure and becomes the product of the old minor by the determinant of a new $3\times 3$ block, namely
\begin{equation*}
\frac{dl_{MR}\wedge dl_{NR}\wedge dl_{PR}}{dx_R\wedge dy_R\wedge dz_R} = \frac{6V_{MNPR}}{l_{MR}\,l_{NR}\,l_{PR}}
%\label{eq_small_thing}
\end{equation*}
(this equality follows from elementary geometry, compare formulas (31) and~(32) in~\cite{JNMP1}). Due to the same considerations as in the case of a $2\to 3$ Pachner move, the minor of $f_3$ gets multiplied by the very same factor given in~(\ref{eq QRQR}). Comparing that with equations~(\ref{eq_tau}) and~(\ref{eq_invariant}), we see that our value~$I(M)$ does not change under a $1\to 4$ Pachner move, as well as under the inverse move.

The fact that no minors considered in our proof vanished obviously implies that the acyclicity of our complex~(\ref{eq_alg_complex}) is preserved under the Pachner moves.
\end{proof}
